% CP-VI-0049
% target_chapter: VI/47 — Čech Cohomology
% source_unit_id: PVI-RECOVER-721195CCB5-L001527
% source: Downloads/theory-of-algebraic-geometry-11.tex L1527-L1752
% solution_provenance: SOURCE_EXPLICIT_SOLUTION

\subsection*{Problem 5: A Plane Curve and Its \v{C}ech Cohomology}

Let \(X\subseteq\mathbb P_k^2\) be the closed subscheme defined by one
homogeneous polynomial
\[
f(x_0,x_1,x_2)=0
\]
of degree \(d\).

Assume
\[
(1,0,0)\notin X.
\]

We first show that \(X\) is covered by the two affine open subsets
\[
U=X\cap D_+(x_1),
\qquad
V=X\cap D_+(x_2).
\]

The complement of
\[
D_+(x_1)\cup D_+(x_2)
\]
inside \(\mathbb P_k^2\) is the single point
\[
(1,0,0).
\]

Since
\[
(1,0,0)\notin X,
\]
every point of \(X\) lies in either \(D_+(x_1)\) or \(D_+(x_2)\).

Hence
\[
X=U\cup V.
\]

Both \(U\) and \(V\) are affine.

On \(D_+(x_1)\), set
\[
t=\frac{x_0}{x_1},
\qquad
z=\frac{x_2}{x_1}.
\]

Then
\[
U\simeq \operatorname{Spec} k[t,z]/(f(t,1,z)).
\]

Similarly, on \(D_+(x_2)\), set
\[
s=\frac{x_0}{x_2},
\qquad
w=\frac{x_1}{x_2}.
\]

Then
\[
V\simeq \operatorname{Spec} k[s,w]/(f(s,w,1)).
\]

Their intersection is
\[
U\cap V=X\cap D_+(x_1x_2).
\]

This is affine.

Since
\[
(1,0,0)\notin X,
\]
we have
\[
f(1,0,0)\neq0.
\]

Therefore the coefficient of \(x_0^d\) in \(f\) is nonzero. After scaling,
we may assume that this coefficient is \(1\).

On the intersection \(U\cap V\), use coordinates
\[
u=\frac{x_0}{x_1},
\qquad
v=\frac{x_2}{x_1}.
\]

Since \(x_2\neq0\) on \(U\cap V\), the coordinate \(v\) is invertible.

Thus
\[
\mathcal O_X(U\cap V)
=
k[u,v,v^{-1}]/(f(u,1,v)).
\]

Because \(f(u,1,v)\) is monic of degree \(d\) in \(u\), every element of
\(\mathcal O_X(U\cap V)\) can be represented uniquely as
\[
\sum_{r=0}^{d-1} a_r(v)u^r,
\qquad
a_r(v)\in k[v,v^{-1}].
\]

The \v{C}ech complex for the cover \(X=U\cup V\) is
\[
0
\longrightarrow
\mathcal O_X(X)
\longrightarrow
\mathcal O_X(U)\oplus\mathcal O_X(V)
\longrightarrow
\mathcal O_X(U\cap V)
\longrightarrow
H^1(X,\mathcal O_X)
\longrightarrow0.
\]

Since \(U\), \(V\), and \(U\cap V\) are affine, this complex computes
\(H^0\) and \(H^1\).

The global sections are the elements of
\[
\mathcal O_X(U)\oplus\mathcal O_X(V)
\]
which agree on \(U\cap V\). For a projective plane curve, the only global
regular functions are constants. Hence
\[
H^0(X,\mathcal O_X)\cong k,
\]
so
\[
\boxed{
	\dim_k H^0(X,\mathcal O_X)=1.
}
\]

Now we compute \(H^1(X,\mathcal O_X)\).

The cokernel of
\[
\mathcal O_X(U)\oplus\mathcal O_X(V)
\longrightarrow
\mathcal O_X(U\cap V)
\]
is represented by the principal parts on the overlap that cannot come from
sections on \(U\) or on \(V\).

A basis for the quotient is given by the monomials
\[
u^r v^{-m}
\]
with
\[
1\leq r,
\qquad
1\leq m,
\qquad
r+m\leq d-1.
\]

Equivalently,
\[
1\leq r\leq d-2,
\qquad
1\leq m\leq d-1-r.
\]

For each \(r=1,\ldots,d-2\), the integer \(m\) can take
\[
d-1-r
\]
values.

Therefore the number of such monomials is
\[
\sum_{r=1}^{d-2}(d-1-r)
=
(d-2)+(d-3)+\cdots+1.
\]

Hence
\[
\dim_k H^1(X,\mathcal O_X)
=
1+2+\cdots+(d-2).
\]

Therefore
\[
\dim_k H^1(X,\mathcal O_X)
=
\frac{(d-1)(d-2)}2.
\]

Thus
\[
\boxed{
	\dim_k H^0(X,\mathcal O_X)=1
}
\]
and
\[
\boxed{
	\dim_k H^1(X,\mathcal O_X)=\frac{(d-1)(d-2)}2.
}
\]

If \(X\) is nonsingular, then
\[
\dim_k H^1(X,\mathcal O_X)
\]
is the genus of the plane curve. Hence a nonsingular plane curve of degree
\(d\) has genus
\[
\boxed{
	g=\frac{(d-1)(d-2)}2.
}
\]
